Las constantes son reales y los signos se mantienen en cada intervalo, salvo las decisiones de prolongación justificadas en cada PVI. Cada expresión se restringe a intervalos donde satisface las condiciones señaladas.
\begin{enumerate}[label=\alph*)]
\item $y=\sqrt{1+3e^{-x^2}}$, \quad $x\in\mathbb R$.
\item $y=x\exp(C/x-1)$, \quad $x>0$.
\item $y=\ln W(x)$, donde $W(x)=Ce^x+\tfrac12(\sin x-\cos x)>0$. Se toma cualquier componente conexa de ese conjunto.
\item $y=0$; o bien $y=\dfrac{1+x^2}{C-x}$, en $(-\infty,C)$ o en $(C,\infty)$.
\item $y+y^3=\dfrac{x^2}{4}+\dfrac{C}{x^2}$, \quad $x>0$. Hay una única raíz real para cada $x$ y $C$.
\item $y=\tan\!\left(\dfrac\pi4e^x\right)-2x$, \quad $x\in(-\infty,\ln2)$.
\item $y=\dfrac{x}{2}+\arctan(Ce^x)+k\pi$, con $k\in\mathbb Z$, y además $y=\dfrac{x}{2}+\dfrac{\pi}{2}+k\pi$. Todas están definidas en $\mathbb R$.
\item $y=x+1$ y $y=x-1$; o bien, con $u=y-x$, $3\ln|u-1|-\ln|u+1|=2x+C$. La familia implícita se restringe a ramas con $u\notin\{-2,-1,1\}$.
\item $y=-x$; o bien, con $v=y/x$, $\dfrac{v^2}{2}-v+2\ln|v+1|=\ln|x|+C$, con $v<-1$ o $v>-1$. Cada rama se define para todo $x>0$ o todo $x<0$.
\item Si $v=(y+1)/(x-2)$, entonces $v+v^3/3=\ln|x-2|+C$. La raíz real $v$ es única; cada solución ocupa uno de los semiejes $x<2$ o $x>2$.
\item $y=x^2\ln(x+C)$, \quad $x>\max\{0,-C\}$.
\item $y=\dfrac{e^x(x+C)}{\sqrt{1-(x+C)^2}}$, \quad $-C-1<x<-C+1$. El numerador conserva el signo de $x+C$.
\item $y=x$ e $y=-x$, ambas en $\mathbb R$. No se pueden concatenar ramas con pendientes opuestas conservando la clase $C^1$.
\item $\ln y+xy=\arctan x+1-\pi/4$, \quad $x>0$, $y>0$. Esta relación determina una única función positiva en $(0,\infty)$.
\item $x=y^2+Ce^y$. Se invierte en cualquier intervalo de $y$ donde $2y+Ce^y\ne0$; su imagen es el intervalo de definición de la rama $y(x)$. No hay soluciones constantes.
\end{enumerate}
